\par\noindent\textbf{(1)} 弦位移在点质量处连续，故t=1+r。点质量垂向受力为T[y\_x(0+)−y\_x(0−)]，令位移随e\textasciicircum{}\{−iωt\}变化，牛顿方程−Mω²t=2iTk r，其中k=ω/c。此式的右侧是左右弦拉力之差，不能错误设成零。
\par\noindent\textbf{(2)} 把Z=T/c=Tk/ω代入跳跃式得−Mω(1+r)=2iZr，所以r=−Mω/(Mω+2iZ)而t=2iZ/(Mω+2iZ)。两者满足t=1+r；给定入射位移归一化后可通过复数代数独立核验两个界面条件。
\par\noindent\textbf{(3)} 左右介质相同，所以R=|r|²=(Mω)²/[(Mω)²+4Z²]、T\_trans=|t|²=4Z²/[(Mω)²+4Z²]。二者相加等于1，点质量在周期内只储存动能、没有净耗散功，因而透射降低并不等价于能量消失。
\par\noindent\textbf{(4)} 低频Mω≪2Z有r≈iMω/(2Z)趋零，质量能跟随弦缓慢运动；高频Mω≫2Z时r→−1、t→0，形成近固定端条件。该无弹簧自由质量没有有限频率内部谐振，若加入弹簧或阻尼才会出现额外谐振与吸收。
